\begin{abstract}
Assuming the continuum hypothesis, we construct an abstract elementary
class with L\"owenheim--Skolem number \(\aleph_0\) that is categorical in
every sufficiently large cardinal but has at least two nonisomorphic
models of cardinality \(\beth_{\omega_2}\).
Thus categoricity does not transfer down to the proposed bound
\(\beth_{(2^{\aleph_0})^+}\), which equals \(\beth_{\omega_2}\) under CH.
Consequently, if ZFC is consistent, the prescribed-threshold
form of Shelah's categoricity conjecture is not provable in ZFC.
\end{abstract}

\section{Introduction}\label{sec:introduction}

Categoricity asks whether a class has exactly one model of a given
cardinality, up to isomorphism. For an abstract elementary class \(K\),
write \(I(K,\kappa)\) for the number of isomorphism types of its models
of cardinality \(\kappa\), and \(\LS(K)\) for its
L\"owenheim--Skolem number. Thus categoricity in \(\kappa\) means
\(I(K,\kappa)=1\). Put
\begin{equation}\label{eq:hanf-function}
 h(\kappa)=\beth_{(2^\kappa)^+},\qquad H(K)=h(\LS(K)).
\end{equation}
Here cardinal successors are identified with their initial ordinals when
used as beth indices, and the beth hierarchy starts at
\(\beth_0=\aleph_0\). All structures and index sets are set-sized, all
languages are finitary, and categoricity includes existence.
The AEC axioms are recalled in Section~\ref{sec:aec}.

\subsection{The prescribed threshold and the result}

For complete first-order theories in a countable language, Morley's
categoricity theorem transfers categoricity from one uncountable
cardinal to every uncountable cardinal \cite{Morley1965}.
Abstract elementary classes retain a coherent notion of strong
substructure and closure under increasing unions while allowing classes
that are not axiomatized by a first-order theory. Shelah introduced
this framework in the mid-1970s for classes of models in infinitary logics
\cite[Section~2]{BoneyVasey2017Survey}.
His categoricity questions for abstract elementary classes have
both qualitative and quantitative forms. Qualitative eventual
categoricity asks for a cardinal \(\chi(\kappa)\), depending only on
the infinite cardinal \(\kappa\), such that, whenever
\(\LS(K)\le\kappa\), categoricity in some
\(\lambda\ge\chi(\kappa)\) implies categoricity in every
\(\mu\ge\chi(\kappa)\). The prescribed-threshold form asks for the
particular bound in~\eqref{eq:hanf-function}:
\begin{equation}\label{eq:target}
 \begin{gathered}
 \text{If \(K\) is categorical in some \(\lambda\ge H(K)\),}\\
 \text{then it is categorical in every \(\mu\ge H(K)\).}
 \end{gathered}
\end{equation}
No amalgamation, joint embedding, tameness, or absence of maximal
models is assumed here. The exact formulation, with both endpoints
included, appears in the discussion of the general conjecture in
\cite[Section~2.3]{SarochTrlifaj2024}.

The number \(H(K)\) is a general existence bound: a model at that size
gives arbitrarily large models by the presentation and
Ehrenfeucht--Mostowski methods; see
\cite[Claim~0.6]{Shelah1999CategoricityAP}.
Existence alone does not assert uniqueness or that every given model
extends to arbitrarily large ones. Categoricity transfers in the
literature obtain stronger conclusions under structural assumptions.
For example, amalgamation and \(\LS(K)\)-tameness yield transfer from a
successor categoricity cardinal at least \(H(K)\) down to \(H(K)\) and
throughout the tail \cite[Corollary~9.10]{Vasey2017Downward}.
The distinction between these assumptions and the bare AEC axioms is
essential for the construction below.

The distinction between a prescribed bound and an unspecified uniform
threshold is already explicit in Grossberg's discussion of the two
forms of the conjecture \cite[Chapter~2, Section~4, Conjectures~4.4--4.6]{Grossberg2021Draft}.
There the possible failure of the prescribed bound is considered
separately from eventual categoricity. A transfer at the prescribed
bound for arbitrary AECs is claimed in
\cite[Theorem~4.1]{Espindola2023Spectra}. The CH construction below has
the opposite conclusion at that bound. Its proof therefore includes
explicit verification of all AEC axioms, as well as separate arguments
for the endpoint models and for the higher categoricity tail.
Section~\ref{sec:booleanization} also identifies a specific obstruction
to the canonical-point step in the cited proof: on the countable-model
site, a dense cosieve need not be lifted by a large categorical model.
The comparison concerns that specific factorization, not alternative
constructions of points.

\begin{theorem}\label{thm:main}
Assume the continuum hypothesis. There is an abstract elementary class
\(K\) in a countable finitary relational language such that
\begin{enumerate}[label=\textup{(\roman*)}]
 \item \(\LS(K)=\aleph_0\);
 \item \(K\) has at least two nonisomorphic models of cardinality
       \(\beth_{\omega_2}=H(K)\);
 \item putting \(\Lambda=\beth_{(2^{\aleph_1})^+}\),
       the class \(K\) is categorical in every cardinal
       \(\mu\ge\Lambda\).
\end{enumerate}
In particular, \(\lambda=\Lambda^+>H(K)\) is a categoricity cardinal
from which downward transfer to \(H(K)\) fails.
\end{theorem}

Theorem~\ref{thm:main} refutes the prescribed-threshold transfer under
CH. Combining the construction with G\"odel's relative consistency of
CH gives the corresponding conditional nonprovability statement;
Corollary~\ref{cor:consistency} states it with a precise first-order
formulation. The same class is categorical on the higher tail beginning
at \(\Lambda\), so the obstruction is compatible with qualitative
eventual categoricity. The construction and both cardinal bounds are
proved here without a qualitative eventual-categoricity theorem.

\subsection{The mechanism of the construction}

There are two principal sorts. The index sort \(I\) controls the size
of the model. The base sort \(B\) is a dense linear order without
endpoints whose bounded intervals are countable; hence
\(|B|\le\aleph_1\). For each nonempty finite ordered tuple of distinct entries from \(I\),
the structure has row fibers indexed by \(\omega\) and column fibers
indexed by \(B\). A row fiber is a torsor for \([B]^{<\omega}\), and a
column fiber for \([\omega]^{<\omega}\), with symmetric difference as
addition: choosing an origin identifies a fiber with its group, but no
origin is named in the language. Choosing one point in each fiber displays
the bit relation as an \(\omega\)-by-\(B\) matrix. Changing row origins
alters finitely many entries in each row, and changing column origins
alters finitely many entries in each column.

When \(B\) is countable, these changes absorb every matrix of bits.
All models in \(K\) with countable base orders and equally sized
index sorts are therefore isomorphic. When \(B\) is uncountable,
the matrix retains a family of colors in \(2^\omega/=^*\), the quotient
of binary sequences by equality outside a finite set. These colors are
meaningful outside a countable set of columns. We impose tests on them,
separately for each tuple. The tests describe countable extensional
relations whose rank labels lie below \(\omega_1\).

In the following overview, \(V_\alpha\) denotes the ordinary cumulative
hierarchy of sets. The technical features can be summarized as follows.
\begin{enumerate}
 \item \emph{Exceptions compatible with strong substructure.}
       An object may have countably many exceptional columns for each
       tuple. A strong inclusion requires every newly added column to
       pass the tests on old tuples. Finite row supports make these
       outside tests exactly independent of old choices of origins.
       This permits both small strong submodels and arbitrary directed
       unions, including smoothness.
 \item \emph{Local rank approximation.}
       An \(\omega_2\)-long sequence of functions
       \(\omega_1\to\omega_1\), increasing modulo countable sets,
       lets countable membership diagrams drawn from
       \(V_{\omega_2}\) pass the rank tests at all but countably many
       coordinates. This produces the model of size
       \(\beth_{\omega_2}\) with uncountable \(B\).
 \item \emph{An obstruction to large index sets.}
       A fixed coherent pattern for every finite arity would give
       arbitrarily large extensional structures with ranks below
       \(\omega_1\), contradicting their collapse into
       \(V_{\omega_1}\). A finite partition argument and an
       unbounded-sample selection produce such a pattern if
       \(|I|\ge\Lambda\) and \(B\) is uncountable.
\end{enumerate}
Finite-support groups, torsors with no named origins, and parity
relations have a close predecessor in the Hart--Shelah examples
\cite[Section~1]{HartShelah1990}. Their canonical structures also use
binary-sequence colors modulo finitely supported sequences,
\(2^\omega/2^{<\omega}\), the same quotient as eventual equality \cite[Section~3]{HartShelah1990}. The present construction
couples finite row and column translations with a locally countable
base order; the surviving colors encode the extensional rank diagrams
needed for the endpoint obstruction. Allowed finite color diagrams and
ranks also play a central role in the coloring classes studied by
Kolesnikov and Lambie-Hanson \cite{KolesnikovLambieHanson2016}.
Here the encoded diagrams are countable, and the full-subsequence tests
supply the exact compatibility used in their limit.

The first mechanism separates the invariance needed for objects from
the exact invariance needed for extensions. The last two mechanisms
compare local approximations, with different exceptional sets, to one
global pattern. Their proofs are given explicitly; in particular, no
infinite homogeneous subset is assumed.

Section~\ref{sec:colors} defines the diagram tests.
Section~\ref{sec:aec} constructs the AEC and verifies all its axioms.
Section~\ref{sec:endpoint} constructs the two endpoint models.
Section~\ref{sec:categoricity} proves the fixed-pattern obstruction,
uses it to bound models with uncountable base, and then establishes
categoricity on the tail. Section~\ref{sec:conclusion} combines the
results and explains the failure of amalgamation and joint embedding.
Section~\ref{sec:booleanization} proves the dense-cosieve obstruction
and makes the precise comparison with the cited topos argument.

\subsection{Conventions}

The construction in Sections~\ref{sec:colors}--\ref{sec:categoricity} is
carried out in ZFC with CH. Throughout those sections,
\begin{equation}\label{eq:standing-cardinals}
 \delta=\omega_1,\qquad \eta=\delta^+=\omega_2,\qquad
 S=2^\omega/=^*,\qquad \xi=2^\delta,\qquad
 \Lambda=\beth_{\xi^+}.
\end{equation}
The relation \(=^*\) means equality except at finitely many natural
numbers. ``Countable'' includes finite, and ``tuple'' in the color
tests always means a nonempty finite ordered tuple with distinct
entries. Subsequences preserve the order of the remaining positions.
The hierarchy \(V_\alpha\) is the ordinary cumulative hierarchy of
sets; it is unrelated to the model sorts. We use choice for set-indexed
families only.
